\section{中心三次元の同変選択と中心商族}
\label{sec:pg-equivariant-central-lifts}
本節では単一のJacobianブロックを扱う。
定階数核の構成を同変に行うことで、固定したactive・中心変数の分解を保つ
還元的群について、古典極限を指定した中心元を同時に選べる。
さらに三変数制限が滑らかな部分では、二次混合項の規格化によって
中心変数による平行移動も処理できる。
平均化そのものは既知のReynolds作用素の性質であり、
ここで用いるのはその定階数核への適用である
\cite[\S2, p.~369]{PDerksenInvariants}。

$W=\CC^3$、$Z=\CC^m$とし、$x=(x,y,z)$と$u$をそれぞれの座標とする。
変形パラメータを一旦$t$と書き、
\eqref{eq:pg-det-quantum}の$\hbar$を$t$に置き換えた代数を$A_{f,t}$とする。
群とその作用を
\begin{equation}
 \begin{aligned}
 G_0&=(\mathrm{GL}(W)\times\mathrm{GL}(Z))\times\mathbf G_m,\\
 G&=\left\{
 \begin{pmatrix}M&C\\0&N\end{pmatrix}
 :M\in\mathrm{GL}(W),\ N\in\mathrm{GL}(Z),\ C\in\operatorname{Hom}(Z,W)
 \right\}\times\mathbf G_m,\\
 X&=Mx+Cu,\qquad U=Nu,\\
 f'(X,U)&=\lambda^{-1}
 f\bigl(M^{-1}(X-CN^{-1}U),N^{-1}U\bigr),\qquad
 t'=\lambda\det(M)t
 \end{aligned}
 \label{eq:pg-equivariant-groups}
\end{equation}
で定める。$G_0\subset G$は$C=0$の部分群である。
代入による同型を$\Phi:A_{f',t'}\to A_{f,t}$とする。
注意\ref{rem:pg-det-block-covariance}により、Weyl同型、割った交換子の核
$\mathcal K$、純中心部分$\mathcal C_m$および商線束
$\mathcal K/\mathcal C_m$は$G$と両立する。
純中心部分への射影$\rho(p)=p(0,u)$は$G_0$と両立する。

\subsection{不変開と中心元の同変持上げ}

\begin{lemma}[アフィン底上の同変持上げ]
\label{lem:pg-equivariant-affine-lifting}
$H$を複素還元的線形代数群、$Y$をアフィン$H$-スキームとする。
次が成り立つ。
\begin{enumerate}
\item $Z_0\subset Y$を$H$-安定閉部分、$V\subset Y$を$Z_0$を含む
$H$-安定開部分とする。
ある$H$-不変関数$q\in\Gamma(Y,\OO_Y)$について、
$q|_{Z_0}=1$かつ$D(q)\subset V$とできる。
\item $\mathcal E$を$Y$上の同変ベクトル束とする。
$\mathcal E|_{Z_0}$の不変切断は$\mathcal E$の不変切断へ持ち上がる。
任意の指標で線形化をtwistした場合も同じである。
\item $Y$上の同変ベクトル束の短完全列は同変に分裂する。
\end{enumerate}
\end{lemma}
\begin{proof}
アフィン底上の同変準連接層の大域切断は有理$H$-加群であり、
各元の軌道の線形包は有限次元である。
その有限次元部分において、Zariski稠密な最大compact部分群で平均すれば
Reynolds射影が得られる。
これは$H$-安定部分加群を保ち、同変線形写像と可換する。

第一の主張では、$Y\setminus V$のイデアルを$I$、$Z_0$のイデアルを$J$とすると
$I+J=\Gamma(Y,\OO_Y)$である。
$p\in I$、$p\equiv1\pmod J$を選び、平均して$q$とすればよい。
第二の主張では、アフィン性による全射
\begin{equation}
 \Gamma(Y,\mathcal E)\longrightarrow
 \Gamma(Z_0,\mathcal E|_{Z_0})
 \label{eq:pg-equivariant-restriction-surjection}
\end{equation}
において任意の持上げを平均する。
第三の主張では、商ベクトル束が有限射影加群に対応するので、
まず$\OO_Y$-線形な分裂を選べる。
Hom加群におけるその平均も$\OO_Y$-線形であり、
合成が恒等写像であるという条件を保つ。
\end{proof}

\begin{theorem}[還元的部分群に関する中心元の同変選択]
\label{thm:pg-equivariant-levi-center}
$\mathcal T_m$を全三次式$f$が滑らかな射影超曲面を定める係数空間とする。
$G_0$-不変な偶関数$Q(f,t)$、$Q(f,0)=1$を用いたprincipal開
\begin{equation}
 \mathcal W_m^0=D(Q)\subset\mathcal T_m\times\mathbb A^1_t
 \label{eq:pg-equivariant-levi-open}
\end{equation}
上で、Weyl係数が正則な中心元$\Omega^0_{f,t}$を、
\begin{equation}
 \begin{aligned}
 \mathcal D\Omega^0_{f,t}&=0,&
 \Omega^0_{f,0}&=f,&
 \rho(\Omega^0_{f,t})&=f_{\mathrm{cent}},\\
 \Omega^0_{f,-t}&=\Omega^0_{f,t},&
 \Phi(\Omega^0_{f',t'})&=\lambda^{-1}\Omega^0_{f,t}\qquad(C=0)
 \end{aligned}
 \label{eq:pg-equivariant-levi-center}
\end{equation}
となるように選べる。
ここで偶性はWeyl係数についての等式である。
さらに$\Omega^0_{f,t}$の$\mathcal K/\mathcal C_m$への像は
この線束の基底となる。
この選択は一般に一意ではない。
\end{theorem}
\begin{proof}
命題\ref{prop:pg-PBW}と
補題\ref{lem:pg-det-divided-kernel}で用いた開集合を、
三変数斉次代数のAS-regular条件、次数三・四のWeyl同型、
および割った交換子の階数$r_m$という条件で定める。
これらの条件は座標変換と両立するので、得られる開集合$\mathcal V_m$は
$G$-安定であり、$\mathcal T_m\times\{0\}$を含む。
$t\mapsto-t$は代数を反対代数へ移し、これらの条件も保つ。

$\mathcal T_m$は判別式を可逆にしたアフィン空間である。
補題\ref{lem:pg-equivariant-affine-lifting}(1)を
$H=G_0\times\boldsymbol\mu_2$と零切断に適用すると、
$\mathcal V_m$内に、不変かつ偶な$Q$による
\eqref{eq:pg-equivariant-levi-open}を取れる。
$t$は$G_0$の指標$\lambda\det(M)$で変換するが、
イデアル$(t)$は安定であるから、$Q(f,0)=1$という条件は平均で保たれる。

このアフィン開上で$\mathcal K$は同変ベクトル束である。
Weyl表示における$\mathcal D$は$t$について偶なので、
$\boldsymbol\mu_2$はWeyl記号を変えずに$\mathcal K$へ作用する。
補題\ref{lem:pg-det-classical-kernel}より、
零切断上の$f$は$\mathcal K|_{t=0}$の切断である。
引戻しの線形化を$\lambda\Phi$へtwistしてこの切断を不変とみなし、
補題\ref{lem:pg-equivariant-affine-lifting}(2)で持ち上げる。
得られた切断を$\widetilde\Omega$とし、
\begin{equation}
 \Omega^0=\widetilde\Omega-\rho(\widetilde\Omega)+f_{\mathrm{cent}}
 \label{eq:pg-equivariant-pure-center-adjustment}
\end{equation}
と置く。
$\rho$は$G_0$および$\boldsymbol\mu_2$と両立し、
純中心三次式は$\mathcal D$で消えるので、
\eqref{eq:pg-equivariant-levi-center}を得る。

商線束における$\Omega^0$の像は零切断上で消えない。
その零点集合を除いた安定開に同補題(1)を再適用する。
必要なら不変な分母を払えば、元のアフィン底でも
$Q(f,0)=1$を満たす不変な偶関数のprincipal開として取れる。
この開では商線束への像が基底となる。
\end{proof}

\begin{theorem}[滑らかな三変数制限上の全ブロック群に関する選択]
\label{thm:pg-equivariant-parabolic-center}
\begin{equation}
 \mathcal T_m^{\mathrm{ell}}
 =\{f\in\mathcal T_m:g(x)=f(x,0)\text{ が滑らかな平面三次を定める}\}
 \label{eq:pg-equivariant-elliptic-base}
\end{equation}
とおく。
$\mathcal T_m^{\mathrm{ell}}\times\mathbb A^1_t$には、零切断全体を含む
$G$-不変なprincipal開$\mathcal W_m^G$と、そこに正則な中心元
$\Omega^G_{f,t}$が存在し、
\begin{equation}
 \mathcal D\Omega^G_{f,t}=0,\qquad
 \Omega^G_{f,0}=f,\qquad
 \Omega^G_{f,-t}=\Omega^G_{f,t},\qquad
 \Phi(\Omega^G_{f',t'})=\lambda^{-1}\Omega^G_{f,t}
 \label{eq:pg-equivariant-parabolic-center}
\end{equation}
を満たす。
商線束$\mathcal K/\mathcal C_m$への像も基底となる。
純中心係数を$f_{\mathrm{cent}}$に固定する条件は課さない。
\end{theorem}
\begin{proof}
滑らかな三変数三次式$g$のアフィン係数空間を$\mathcal T_0$とする。
方向微分の写像
\begin{equation}
 \delta_g:W\longrightarrow\operatorname{Sym}^2W^*,
 \qquad v\longmapsto\partial_vg
 \label{eq:pg-equivariant-gradient-embedding}
\end{equation}
は各点で単射である。
実際、非零方向で$\partial_vg=0$なら、$g$はその方向の変数に依存せず、
その座標点が平面三次の特異点となる。
従って余核もベクトル束である。
二次形式側の作用に$\lambda^{-1}$を含めれば
\eqref{eq:pg-equivariant-gradient-embedding}は
$\mathrm{GL}(W)\times\mathbf G_m$-同変である。
補題\ref{lem:pg-equivariant-affine-lifting}(3)によって、
正則な同変左逆$\sigma_g$を一つ選ぶ：
\begin{equation}
 \sigma_g:\operatorname{Sym}^2W^*\longrightarrow W,
 \qquad \sigma_g\delta_g=\mathrm{id}_W.
 \label{eq:pg-equivariant-gradient-retraction}
\end{equation}

$f$のactive次数二の部分を$q_f(x,u)=\sum_\alpha q_\alpha(x)u_\alpha$と書き、
\begin{equation}
 B_f(u)=\sum_\alpha\sigma_g(q_\alpha)u_\alpha,
 \qquad
 f^0(X,u)=f(X-B_f(u),u)
 \label{eq:pg-equivariant-centering}
\end{equation}
と定める。
全ての係数は$\mathcal T_m^{\mathrm{ell}}$上で正則であり、
$f^0$の二次混合部分には$\sigma_g$を適用すると零になる。
$G$の作用下での変換を計算すると、
\begin{equation}
 B_{f'}=(MB_f-C)N^{-1},\qquad
 X+B_{f'}U=M(x+B_fu),\qquad
 (f')^0(X,U)=\lambda^{-1}f^0(M^{-1}X,N^{-1}U)
 \label{eq:pg-equivariant-centering-covariance}
\end{equation}
となる。
第一の等式では、$g$を$x\mapsto x-Cu$で移すと
二次混合項が$-\delta_g(Cu)$だけ変わることと
$\sigma_g\delta_g=\mathrm{id}$を用いた。
従って$(f,t)\mapsto(f^0,t)$は中心平行移動に関して不変で、
残る群$G_0$について同変である。

定理\ref{thm:pg-equivariant-levi-center}の開と中心元を用いて
\begin{equation}
 \begin{aligned}
 \mathcal W_m^G
 &=\{(f,t):(f^0,t)\in\mathcal W_m^0\},\\
 \Omega^G_{f,t}(x,u)
 &=\Omega^0_{f^0,t}(x+B_f(u),u)
 \end{aligned}
 \label{eq:pg-equivariant-centered-quantization}
\end{equation}
と置く。
開集合は不変関数$Q(f^0,t)$で定まり、$Q(f^0,0)=1$である。
中心変数による線形代入はWeyl対称化と可換し、代数の同型を与えるので、
中心性、偶性および商線束の生成性が保たれる。
\eqref{eq:pg-equivariant-centering-covariance}と
\eqref{eq:pg-equivariant-levi-center}から全$G$に関する共変性を得る。
古典極限は$f^0(x+B_fu,u)=f(x,u)$である。
\end{proof}

\subsection{中心商の族と規格化の条件}

\begin{corollary}[同変な正則中心商族]
\label{cor:pg-equivariant-central-quotient}
前二定理のそれぞれの底と群について、
$\mathcal J=\OO\Omega\subset\mathcal A_3$は同変な中心直線部分束である。
各幾何学的ファイバーで$\Omega$は左右正則であり、
中心商$\mathcal B=\mathcal A/(\mathcal J)$の各次数部分は底上局所自由である。
そのファイバーは
\begin{equation}
 \operatorname{Hilb}_{\mathcal B}(z)
 =\frac{1-z^3}{(1-z)^{m+3}}
 \label{eq:pg-equivariant-quotient-hilbert}
\end{equation}
を持ち、左右injective dimension $m+2$、AS parameter $m$の
AS-Gorenstein代数である。

変形パラメータを群作用から独立に固定するには、
\begin{equation}
 t=\eta\hbar,\qquad
 \eta'=\lambda\det(M)\eta,\qquad \hbar'=\hbar
 \label{eq:pg-equivariant-fixed-parameter}
\end{equation}
とおく。
この代入による引戻しは、それぞれ
\begin{equation}
 \left[(\mathcal T_m\times\mathbb A^1_\eta)/G_0\right]
       \times\mathbb A^1_\hbar,
 \qquad
 \left[(\mathcal T_m^{\mathrm{ell}}\times\mathbb A^1_\eta)/G\right]
       \times\mathbb A^1_\hbar
 \label{eq:pg-equivariant-quotient-stacks}
\end{equation}
の、$\hbar=0$全体を含む開部分上に次数付き代数の族として降りる。
特殊ファイバーは$\CC[x,u]/(f)$であり、一次括弧は
$\eta\{-,-\}_f$である。
\end{corollary}
\begin{proof}
共変性は$\mathcal J$およびその生成する両側イデアルを保つ。
命題\ref{prop:pg-PBW}により、各ファイバーで
\begin{equation}
 \operatorname{gr}_x A_{f,t}
 \simeq S_{g,t}\otimes_\CC\CC[u_1,\ldots,u_m]
 \label{eq:pg-equivariant-fiber-domain}
\end{equation}
はdomainであるから、$A_{f,t}$もdomainである。
商線束を生成する$\Omega$は零ではなく、従って左右正則である。
各次数の乗法写像
\begin{equation}
 \mathcal J\otimes\mathcal A_{d-3}\longrightarrow\mathcal A_d
 \label{eq:pg-equivariant-regular-multiplication}
\end{equation}
は全ファイバーで単射なので、ベクトル束の写像として局所的に分裂する。
その余核$\mathcal B_d$は局所自由であり、Hilbert級数も従う。
$u_1,\ldots,u_m$はPBW性による中心正則列で、
その商$S_{g,t}$は三次元Koszul AS-regularである。
補題\ref{lem:pg-central-linear}を$m$回適用し、
さらに補題\ref{lem:pg-central-cubic}を適用するとAS条件を得る。

\eqref{eq:pg-equivariant-fixed-parameter}は群作用と両立する。
元の開は$t=0$全体を含むので、引戻した開は$\hbar=0$全体を含む。
同変代数と同変イデアルの降下により
\eqref{eq:pg-equivariant-quotient-stacks}上の族が得られる。
その平坦性はアフィンatlas上で示した次数別局所自由性から従う。
定義関係式を$\hbar$で割って特殊化すれば一次括弧は
$\eta\{-,-\}_f$となる。
\end{proof}

\begin{remark}[純中心部分がない単一ブロック]
\label{rem:pg-equivariant-line-and-gerbe}
$m=0$では$\mathcal C_m=0$なので、定階数核$\mathcal K$自身が
正準な同変中心直線である。
従って同変な中心イデアル$(\mathcal K)$とその商には、
中心元の基底やReynolds平均の選択は不要である。
これは単一ブロックについての主張であり、複数ブロックの中心元の和の
相対規格化を定めるものではない。

座標スカラーの作用と射影座標のモジュライへの降下は、後の命題
\ref{prop:pg-equivariant-poisson-gerbe}と系
\ref{cor:pg-equivariant-character-descent}で区別して扱う。
\end{remark}

\begin{remark}[純中心係数の固定と平行移動]
\label{rem:pg-equivariant-pure-center-obstruction}
$m\geq1$では、零切断全体を含む$G$-不変開上で、全$G$に関する共変性と
$\rho(\Omega_f)=f_{\mathrm{cent}}$を同時に要求することは一般にできない。
実際、中心平行移動$X=x+Cu$とこの規格化を全ての$C$に適用すると、
$\Omega_f$のWeyl記号$P_f(x,u)$について
\begin{equation}
 P_f(-Cu,u)=f(-Cu,u)\qquad\text{for all }C
 \label{eq:pg-equivariant-pure-center-evaluation}
\end{equation}
を得る。
$u\ne0$なら$Cu$は任意の$x\in W$を動くので、$P_f=f$が強制される。
しかし例えば
\begin{equation}
 f=\frac{x^3+y^3+z^3}{3}+\sum_{\alpha=1}^m u_\alpha^3
 \label{eq:pg-equivariant-pure-center-fermat}
\end{equation}
について、$t\ne0$の適切な開上では$\sym(f)$は中心元ではない。
定理\ref{thm:explicit-hesse-center}のFermat公式が与える商線束の基底は
$(x^3+y^3+z^3)/3-2t^2\sym(xyz)/3$であり、
Weyl基底の独立性から$(x^3+y^3+z^3)/3$はこの直線に入らない。
従って定理\ref{thm:pg-equivariant-parabolic-center}で
純中心係数の固定を課さないことは必要である。
同定理は滑らかな三変数制限の上で全$G$を扱う。
三変数制限が特異または零の場合を含む全$\mathcal T_m$については、
定理\ref{thm:pg-equivariant-levi-center}の$G_0$-同変選択までを結論している。
\end{remark}

\subsection{階数二のPoissonモジュライへの降下}

\begin{proposition}[階数二のPoissonスタックと座標スカラーのgerbe]
\label{prop:pg-equivariant-poisson-gerbe}
$m\ge2$、$n=m+1$、$V=W\oplus Z$、$\dim W=3$、$\dim Z=m$とする。
\eqref{eq:pg-equivariant-elliptic-base}の係数空間と
\eqref{eq:pg-equivariant-groups}の群$G$について、
\begin{equation}
 \begin{aligned}
 \mathscr S_m&=\mathcal T_m^{\mathrm{ell}}\times\mathbf G_{m,\eta},\\
 K_{\mathrm{sc}}&=
 \{(aI_3,aI_m,0,a^{-3}):a\in\mathbf G_m\}\subset G,
 \qquad \overline G=G/K_{\mathrm{sc}}
 \end{aligned}
 \label{eq:pg-equivariant-gerbe-parameter}
\end{equation}
とおく。$G$は$f$に\eqref{eq:pg-equivariant-groups}で、
$\eta$に$\eta'=\lambda\det(M)\eta$で作用する。

$\mathfrak P_{n,1}^{\mathrm{ell}}$を、相対Poissonスタック
$\mathfrak{Pois}^{\mathrm{rel}}_n$の一般階数二の
\emph{被約な局所閉層}のうち、表示$\pi=\iota_{dF}P$の
support三次元空間$W_P$への制限$F|_{W_P}$が
滑らかな平面三次を定める非空稠密開部分とする。
このとき
\begin{equation}
 [\mathscr S_m/\overline G]\simeq
 \mathfrak P_{n,1}^{\mathrm{ell}},\qquad
 [\mathscr S_m/G]\longrightarrow
 \mathfrak P_{n,1}^{\mathrm{ell}}
 \quad\text{は}\ \mathbf G_m\text{-gerbe}
 \label{eq:pg-equivariant-poisson-gerbe}
\end{equation}
である。ここで開部分という条件は階数二の層の中での条件であり、
一般には相対Poissonスタック全体の開部分ではない。
特に$n=3$では底スタックの次元は
$\dim\mathfrak P_{3,1}^{\mathrm{ell}}=17$であり、
$[\mathscr S_2/G]$の次元は$16$である。
\end{proposition}
\begin{proof}
$\mathsf P_W\subset\mathrm{GL}(V)$を$W$を保つ放物部分群とする。
従って$G=\mathsf P_W\times\mathbf G_{m,\lambda}$である。
$P_0=e_1\wedge e_2\wedge e_3\in\bigwedge^3W$を固定し、
\begin{equation}
 \widetilde{\mathscr R}^{\mathrm{ell}}
 =\{(F,P):V(F)\text{は滑らか},\quad
 P\ne0\text{は分解可能},\quad
 F|_{\operatorname{supp}P}\text{は滑らかな平面三次を定める}\}
 \label{eq:pg-equivariant-decomposable-incidence}
\end{equation}
とおく。これは滑らかな三次式のアフィン係数空間と
非零Pl\"ucker錐の積の開部分である。
非零の分解可能三重ベクトルはsupportの三次元空間と
その最高外積の非零元を一意に定める。
$\mathrm{GL}(V)$は三次元部分空間全体に推移的に作用するので、
\begin{equation}
 \mathrm{GL}(V)\times^{\mathsf P_W}\mathscr S_m
 \xrightarrow{\ \sim\ }\widetilde{\mathscr R}^{\mathrm{ell}},\qquad
 [g,(f,\eta)]\longmapsto
 \bigl(f\circ g^{-1},\eta(\bigwedge^3g)P_0\bigr)
 \label{eq:pg-equivariant-associated-incidence}
\end{equation}
を得る。ここでは$\mathsf P_W$の作用は$\lambda=1$の作用である。
この同型は任意の底の上でも成立する。
実際、supportの部分束に適合する枠を局所的に選べば逆写像を構成でき、
枠の違いはちょうど$\mathsf P_W$の作用となる。

階数層のスキーム構造についても確認する。
縮約写像$\bigwedge^2V^*\to V$が階数三である局所閉条件を課すと、
その像は階数三の部分束$\mathcal W\subset V\otimes\OO$となる。
全ての二重縮約が$\mathcal W$に入ることから、
$P\in\bigwedge^3\mathcal W$である。
像の階数が三なので、この最高外積の係数は局所的に単元である。
従ってこの行列式的な階数三の層は、非零Pl\"ucker錐と
スキームとして同型である。
その上では$\iota_{dF}P$は三変数Jacobian括弧であり、
任意の$F$に対してJacobi恒等式を満たす。
よってJacobi方程式は追加条件を課さない。
正規形定理による一般階数二との対応と合わせると、
\eqref{eq:pg-equivariant-decomposable-incidence}は
本稿の被約な階数二の層のアフィン三次式による表示である。
全Jacobiスキームの他の層やその厚みについての同定は用いない。
滑らかな三次超曲面の一般の平面切断はBertiniの定理により滑らかなので、
この開条件は各滑らかな三次式の上で非空である。
非零Pl\"ucker錐と滑らかな三次式の係数空間は既約であり、
その開部分の商である上記の層でも稠密となる。

三次式の代表を忘れる自由な作用は
\begin{equation}
 \lambda\cdot(F,P)=(\lambda^{-1}F,\lambda P)
 \label{eq:pg-equivariant-equation-scalar}
\end{equation}
である。その商を$\mathscr R^{\mathrm{ell}}$と書けば、
$\mathscr R^{\mathrm{ell}}$は
$\bigwedge^3V\otimes\OO_{U_N}(-1)$内の上記の局所閉部分であり、
$\mathfrak P_{n,1}^{\mathrm{ell}}
 =[\mathscr R^{\mathrm{ell}}/\mathrm{PGL}(V)]$である。
一方、\eqref{eq:pg-equivariant-associated-incidence}から
\begin{equation}
 [\mathscr S_m/G]
 \simeq
 [\widetilde{\mathscr R}^{\mathrm{ell}}/
       (\mathrm{GL}(V)\times\mathbf G_{m,\lambda})]
 \label{eq:pg-equivariant-two-quotients}
\end{equation}
を得る。右辺の作用では
$(aI_V,a^{-3})$が$(F,P)$を固定する。
この二つのスカラーを区別するため、
\begin{equation}
 \widehat G=(\mathrm{GL}(V)\times\mathbf G_{m,\lambda})/K_{\mathrm{sc}},
 \qquad
 1\longrightarrow\mathbf G_{m,\lambda}
 \longrightarrow\widehat G
 \longrightarrow\mathrm{PGL}(V)\longrightarrow1
 \label{eq:pg-equivariant-effective-group-sequence}
\end{equation}
と書く。左の単射は$\lambda\mapsto[(I_V,\lambda)]$である。
実際、fppf局所的に$\mathrm{PGL}(V)$への射の核の元を$[(aI_V,\lambda)]$と表せば、
$[(I_V,a^3\lambda)]$と一意に書ける。
従ってまず$K_{\mathrm{sc}}$で割り、次に自由な
\eqref{eq:pg-equivariant-equation-scalar}の作用で割ると、
$\mathrm{PGL}(V)$による商を得る。
これが\eqref{eq:pg-equivariant-poisson-gerbe}の第一の同型を与える。
$K_{\mathrm{sc}}$は$G$の中心にあり$\mathscr S_m$に自明に作用するので、
第二の射はこの$K_{\mathrm{sc}}\simeq\mathbf G_m$をbandとするgerbeである
（商のgerbe判定について\cite[Lemma~78.27.1, Tag~\texttt{06PD}]{PStacks}も参照）。
特に幾何学的点$(f,\eta)$における安定化群の関係は、fppf層としての完全列
\begin{equation}
 1\longrightarrow K_{\mathrm{sc}}
 \longrightarrow G_{(f,\eta)}
 \longrightarrow\operatorname{Aut}(X,\pi)\longrightarrow1
 \label{eq:pg-equivariant-stabilizer-sequence}
\end{equation}
である。右端は射影的なPoisson対の自己同型群であり、
方程式の代表を変える自由なスカラーを安定化群に余分に含めていない。

$n=3$、$m=2$では、三次式の係数空間の次元は$35$なので、
\begin{equation}
 \dim\mathfrak P_{3,1}^{\mathrm{ell}}
 =35+1-(9+4+6+1)+1=17.
 \label{eq:pg-equivariant-threefold-stack-dimension}
\end{equation}
$\overline G$の代わりに$G$で割るgerbeの次元は一小さい。
\end{proof}

\begin{corollary}[指標による次数線形化の降下]
\label{cor:pg-equivariant-character-descent}
$\gcd(m,3)=1$とする。整数$r,s$を$mr-3s=1$となるように選び、
\begin{equation}
 \theta:G\longrightarrow\mathbf G_m,
 \qquad \theta(M,N,C,\lambda)=(\det N)^r\lambda^s,
 \qquad H=\ker\theta
 \label{eq:pg-equivariant-scalar-character}
\end{equation}
とおく。この選択により$G\simeq K_{\mathrm{sc}}\times H$、
$H\simeq\overline G$となり、
\eqref{eq:pg-equivariant-poisson-gerbe}のgerbeは中立になる。

系\ref{cor:pg-equivariant-central-quotient}の族を
$\eta\ne0$へ制限する。各次数$d$の線形化を指標$\theta^{-d}$で
twistすると、周囲代数、中心三次イデアルおよびその商は、
\begin{equation}
 \mathcal U\subset
 \mathfrak P_{n,1}^{\mathrm{ell}}\times\mathbb A^1_\hbar,
 \qquad
 \mathfrak P_{n,1}^{\mathrm{ell}}\times\{0\}\subset\mathcal U
 \label{eq:pg-equivariant-rigidified-family-base}
\end{equation}
上の次数付き$\OO$-代数の族へ降りる。
その各次数部分は局所自由で、幾何学的ファイバーのHilbert級数、
AS-Gorenstein性、古典極限および一次Poisson括弧は
系\ref{cor:pg-equivariant-central-quotient}のものと同じである。
この族の選択は、先の中心元の選択と指標$\theta$に依存してよい。
特に$m=2$、$n=3$では
\begin{equation}
 \theta=(\det N)^2\lambda
 \label{eq:pg-equivariant-threefold-character}
\end{equation}
を取れるため、上の十七次元の底スタック上でこの代数族を得る。
\end{corollary}
\begin{proof}
$k(a)=(aI_3,aI_m,0,a^{-3})$と書くと、
$\theta(k(a))=a^{mr-3s}=a$である。
$K_{\mathrm{sc}}$は中心部分群なので、
\begin{equation}
 K_{\mathrm{sc}}\times H\xrightarrow{\ \sim\ }G,
 \qquad (k,h)\longmapsto kh,
 \qquad
 g\longmapsto
 \bigl(k(\theta(g)),k(\theta(g))^{-1}g\bigr)
 \label{eq:pg-equivariant-character-splitting}
\end{equation}
は群スキームの同型である。
$K_{\mathrm{sc}}$が底に自明に作用することから、
$[\mathscr S_m/G]\simeq[\mathscr S_m/H]\times B\mathbf G_m$であり、
$[\mathscr S_m/H]\simeq\mathfrak P_{n,1}^{\mathrm{ell}}$を得る。

元の次数$d$の線形化を$\rho_d(g)$と書き、
\begin{equation}
 \rho_d^{\theta}(g)=\theta(g)^{-d}\rho_d(g)
 \label{eq:pg-equivariant-character-twist}
\end{equation}
へ換える。
$K_{\mathrm{sc}}$の元$k(a)$は元の代数の次数$d$部分に$a^d$で作用するので、
新しい線形化では自明に作用する。
従って$\rho_d^{\theta}$は$\overline G$を通じて作用する。
積は次数を加えるため
$\theta^{-d}\theta^{-e}=\theta^{-(d+e)}$となり、積と単位の同変性は保たれる。
中心三次直線の線形化も$\theta^{-3}$でtwistすれば、
その包含および生成するイデアルの乗法写像は同変である。
従って代数と商は$\overline G$による商へ降りる。
ここでは商スタック上の準連接層と同変準連接層の同値
\cite[Proposition~96.14.3, Tag~\texttt{06WT}]{PStacks}を
次数別に用いており、積とイデアルも同変写像として降りる。

底の開集合は$G$-不変であり$K_{\mathrm{sc}}$が底に自明に作用するので、
\eqref{eq:pg-equivariant-rigidified-family-base}へ降りる。
線形化の変更はアフィンatlas上の代数、中心イデアル、乗法を変えない。
従って次数別局所自由性と各ファイバーの主張は前の系から従う。
$m=2$では$r=2,s=1$が所要の等式を満たす。
\end{proof}

\begin{remark}[降下した古典極限と中心三次直線]
\label{rem:pg-equivariant-classical-family}
前の系の降下した代数を$\mathcal A^\theta$、$\mathcal B^\theta$とし、
$E_\theta=\mathcal A^\theta_1|_{\hbar=0}$を次数一のベクトル束とする。
古典極限は底スタック上で
\begin{equation}
 \mathcal A^\theta|_{\hbar=0}\simeq\operatorname{Sym}(E_\theta),
 \qquad
 \mathcal B^\theta|_{\hbar=0}
 \simeq\operatorname{Sym}(E_\theta)/(\mathcal J_0)
 \label{eq:pg-equivariant-classical-family}
\end{equation}
と書ける。ここで$\mathcal J_0\subset\operatorname{Sym}^3(E_\theta)$は、
局所的に古典方程式$f$が張る中心三次直線であり、
右辺ではそれが生成する斉次イデアルを取る。
この表示は各atlasでの多項式環と方程式の表示を降ろしたものである。
指標による一様な次数twistは相対$\operatorname{Proj}$を変えないので、
その射影化は底Poissonスタック上の元の普遍三次超曲面族となる。
降下するのは中心三次直線とその生成イデアルであり、
底スタック上でこの直線を自明化する一つの大域生成元$\Omega$を
選べることまでは主張しない。
\end{remark}

\begin{remark}[次数重みと十七次元の意味]
\label{rem:pg-equivariant-gerbe-degree-weights}
\eqref{eq:pg-equivariant-fixed-parameter}によって$\hbar$を固定したときも、
$K_{\mathrm{sc}}$は$(f,\eta,\hbar)$に自明に作用する。
しかし座標代入は各次数一生成元を$a$倍するので、
周囲代数$\mathcal A_d$および中心商$\mathcal B_d$には$a^d$で作用する。
中心三次直線には$a^3$で作用するため、そのイデアルは同変であるが、
次数付き代数のこの非自明な重みは消えない。
従って元の線形化のままの代数はgerbe上の族である。
前の系は$\gcd(m,3)=1$のとき、この次数重みを指標によって取り除いている。
$3\mid m$なら、$G$の任意の指標の$K_{\mathrm{sc}}$上の重みは$3$の倍数である。
実際、指標は$\det M,\det N,\lambda$の整数冪の積であり、
それらの重みはそれぞれ$3,m,-3$である。
この場合は上の指標による分裂は得られないが、
gerbeが非中立であることや別の降下法の不存在までは結論しない。

十七という数は$n=3$での底Poissonスタックの次元である。
固定した$\hbar\ne0$における抽象的な非可換代数のモジュライへの写像が
generically finiteであることや、その像の次元が十七であることは主張しない。
一方、$\hbar$を保つ次数付き変形族全体の同型は、
特殊ファイバーの可換代数と$\hbar^{-1}[-,-]\bmod\hbar$を保つ。
従ってその射影化から元のPoisson対$(X,\pi)$が回復される。
これは一般ファイバーだけからの回復とは異なる。
\end{remark}
